% Source: https://artofproblemsolving.com/community/c3h2673134_a_proving_problem
% Stable identity: aops-topic-2673134-q01
By the binomial theorem, $(m+1)^{k+1}-m^{k+1}=\displaystyle\sum_{j=0}^{k}\dbinom{k+1}{j}m^j$. Summing over $m=1,\dots,n$, the left side telescopes:
\[(n+1)^{k+1}-1=\sum_{j=0}^{k}\dbinom{k+1}{j}S_j(n).\]
The $j=k$ term is $(k+1)S_k(n)$, which gives the formula.

\emph{Polynomial.} By induction on $k$: if each $S_j$ with $j<k$ is a polynomial of degree $j+1\le k$, the right side is a polynomial whose only term of degree $k+1$ is $n^{k+1}$ (from $(n+1)^{k+1}$). So $S_k(n)$ is a polynomial of degree $k+1$ with leading coefficient $\dfrac1{k+1}$. (The base case $S_0(n)=n$ holds.)

\emph{The case $k=3$.} With $S_1=\dfrac{n(n+1)}2$ and $S_2=\dfrac{n(n+1)(2n+1)}6$,
\[4S_3=(n+1)^4-1-n-4S_1-6S_2=(n+1)^4-(n+1)-2n(n+1)-n(n+1)(2n+1).\]
Factoring out $n+1$: $(n+1)\left[(n+1)^3-1-2n-2n^2-n\right]=(n+1)\left(n^3+n^2\right)=n^2(n+1)^2$. So
\[S_3(n)=\boxed{\left(\dfrac{n(n+1)}{2}\right)^2}.\]
